\medskip

\emph{Pour chacune des affirmations suivantes, indiquer si elle est vraie ou fausse.\\
Chaque réponse doit être justifiée. Une réponse non justifiée ne rapporte aucun point.
}

\medskip

\begin{enumerate}
\item On considère la fonction $f$ définie sur $\R$ par 
\[f(x) = \e^x + x.\]

\textbf{Affirmation A} : La fonction $f$ admet pour tableau de variations le tableau ci- dessous:

\begin{center}
\psset{unit=1cm,arrowsize=2pt 3}
\begin{pspicture}(12,2.5)
\psframe(12,2.5)\psline(0,2)(12,2)\psline(3,0)(3,2.5)
\uput[u](1.5,1.9){$x$} \uput[u](3.4,1.9){$- \infty$} \uput[u](11.6,1.9){$+\infty$}
\uput[u](3.4,0){$- \infty$}\uput[d](11.6,2){$+ \infty$} \rput(1.5,1){ variations de $f$}
\psline{->}(3.8,0.3)(11.2,1.7)
\end{pspicture}
\end{center}

\textbf{Affirmation B} : L'équation $f(x) = -2$ admet deux solutions dans $\R$.
\item \textbf{Affirmation C} :

\[\displaystyle\lim_{x \to + \infty} \dfrac{\ln(x) - x^2 +2}{3x^2} = - \dfrac13.\]

\item On considère la fonction $k$ définie et continue sur $\R$ par

\[k(x) = 1 + 2\e^{-x^2 + 1}.\]

\textbf{Affirmation D} : Il existe une primitive de la fonction $k$ décroissante sur $\R$.

\item On considère l'équation différentielle

\[(E): \quad 3y' + y = 1.\]

\textbf{Affirmation E} : La fonction $g$ définie sur $\R$ par 
\[g(x) = 4\e^{- \frac13 x} + 1\]

est solution de l'équation différentielle $(E)$ avec $g(0) = 5$.
\item \textbf{Affirmation F} : Une intégration par parties permet d'obtenir : 

\[\displaystyle\int_0^1 x\e^{-x}\:\text{d}x = 1 - 2\e^{-1}.\]

\end{enumerate}

\bigskip
